\chapter{Marco teórico financiero}
\label{ch:marco-teorico}

Este capítulo presenta los fundamentos teóricos sobre los que se construye
este trabajo. Se comienza con los conceptos de rentabilidad y riesgo, se
desarrolla el modelo de selección de carteras de Markowitz y se introducen
el ratio de Sharpe y el modelo CAPM como extensiones naturales.
La notación que se establece aquí se mantendrá consistente a lo largo de
toda la memoria.

%--------------------------------------------------------------------
\section{Rentabilidad, riesgo y covarianza}
%--------------------------------------------------------------------

\subsection{Rentabilidad de un activo}

Sea $P_t$ el precio ajustado de un activo financiero en el instante $t$.
La \textbf{rentabilidad simple} entre $t-1$ y $t$ se define como:
\begin{equation}
    R_t = \frac{P_t - P_{t-1}}{P_{t-1}} = \frac{P_t}{P_{t-1}} - 1.
    \label{eq:rentabilidad-simple}
\end{equation}

En este trabajo se emplea la \textbf{rentabilidad logarítmica} (o
log-retorno), definida como:
\begin{equation}
    r_t = \ln \frac{P_t}{P_{t-1}} = \ln(P_t) - \ln(P_{t-1}).
    \label{eq:rentabilidad-log}
\end{equation}

La rentabilidad logarítmica tiene dos ventajas para el análisis cuantitativo.
Es \textbf{aditiva en el tiempo}: la rentabilidad acumulada en $k$ días es la
suma de las diarias,
\begin{equation}
    r_{t \to t+k} = \ln \frac{P_{t+k}}{P_t} = \sum_{j=1}^{k} r_{t+j},
    \label{eq:aditividad-log}
\end{equation}
lo que permite agregar a frecuencias mayores con una simple suma ---así se pasa
de datos diarios a la frecuencia mensual del backtesting---. Y para variaciones
pequeñas ($|R_t| \ll 1$) cumple $r_t \approx R_t$, con lo que hereda la
interpretabilidad de la rentabilidad simple.

\subsection{Riesgo: volatilidad}

El \textbf{riesgo} de un activo se mide con la \textbf{volatilidad}
$\sigma$, definida como la desviación típica de sus rentabilidades:
\begin{equation}
    \sigma = \sqrt{\Var[r]} = \sqrt{\frac{1}{T-1} \sum_{t=1}^{T}
    \bigl(r_t - \bar{r}\bigr)^2},
    \label{eq:volatilidad}
\end{equation}
donde $\bar{r} = \frac{1}{T}\sum_{t=1}^T r_t$ es la rentabilidad media
muestral. La volatilidad se expresa típicamente en términos anualizados;
si $r$ es una rentabilidad diaria, $\sigma_{\text{anual}} = \sigma \cdot
\sqrt{252}$, y si $r$ es mensual, $\sigma_{\text{anual}} = \sigma \cdot
\sqrt{12}$.

\subsection{Covarianza y correlación}

Dados dos activos $i$ y $j$, su \textbf{covarianza} mide la tendencia
conjunta de sus rentabilidades:
\begin{equation}
    \sigma_{ij} = \Cov[r_i, r_j]
    = \frac{1}{T-1} \sum_{t=1}^{T} (r_{it} - \bar{r}_i)(r_{jt} - \bar{r}_j).
    \label{eq:covarianza}
\end{equation}

La \textbf{matriz de covarianza} $\boldsymbol{\Sigma} = [\sigma_{ij}] \in
\mathbb{R}^{n \times n}$ recoge las covarianzas entre todos los pares de
activos de un universo de $n$ títulos. Es simétrica ($\sigma_{ij} =
\sigma_{ji}$) y semidefinida positiva ($\mathbf{x}^T \boldsymbol{\Sigma}
\mathbf{x} \ge 0$ para todo $\mathbf{x} \in \mathbb{R}^n$).

La \textbf{correlación} normaliza la covarianza al intervalo $[-1, 1]$:
\begin{equation}
    \rho_{ij} = \frac{\sigma_{ij}}{\sigma_i \, \sigma_j},
    \label{eq:correlacion}
\end{equation}
donde $\sigma_i = \sqrt{\sigma_{ii}}$ es la volatilidad del activo $i$.

\subsection{El principio de diversificación}

La varianza de una cartera con pesos $\mathbf{w} = (w_1, \dots, w_n)^T$
es:
\begin{equation}
    \sigma_p^2 = \mathbf{w}^T \boldsymbol{\Sigma} \mathbf{w}
    = \sum_{i=1}^{n} w_i^2 \sigma_i^2
    + \sum_{i=1}^{n} \sum_{\substack{j=1 \\ j \neq i}}^{n} w_i w_j \sigma_{ij}.
    \label{eq:varianza-cartera}
\end{equation}

El primer sumando recoge las \textbf{varianzas propias} de cada activo; el
segundo, los \textbf{términos de covarianza} entre pares de activos. Es este
segundo bloque el que hace que el riesgo de la cartera no sea, en general, la
suma ponderada de los riesgos individuales.

En efecto, cuando los activos no están perfectamente correlacionados, la
desviación típica de la cartera satisface
\begin{equation}
    \sigma_p \le \sum_{i=1}^{n} w_i \sigma_i \qquad (w_i \ge 0),
    \label{eq:subaditividad}
\end{equation}
con igualdad si y solo si todas las correlaciones valen $1$. Combinar activos
imperfectamente correlacionados reduce, por tanto, el riesgo por debajo de la
media ponderada de los individuales sin sacrificar rentabilidad esperada: esta
es la esencia de la \textbf{diversificación} y el punto de partida de Markowitz.
Su descomposición en riesgo \emph{diversificable} (idiosincrásico) y \emph{no
diversificable} (sistemático) se formaliza con modelos de factores como el CAPM
y el APT (Capítulo~\ref{ch:limitaciones}).

%--------------------------------------------------------------------
\section{Teoría Moderna de Carteras de Markowitz}
%--------------------------------------------------------------------

En su artículo seminal de 1952, Harry Markowitz propuso un marco
matemático para construir carteras óptimas a partir de dos criterios
contrapuestos: maximizar la rentabilidad esperada y minimizar el riesgo
\cite{Markowitz1952}.

\subsection{Rentabilidad y riesgo de una cartera}

Sea $\mathbf{w} \in \mathbb{R}^n$ el vector de pesos de una cartera,
donde $w_i$ representa la fracción de capital invertida en el activo $i$.
Sea $\boldsymbol{\mu} = (\mu_1, \dots, \mu_n)^T$ el vector de
rentabilidades esperadas de los activos, con $\mu_i = \E[r_i]$.

La \textbf{rentabilidad esperada} de la cartera es la combinación lineal
de las rentabilidades esperadas de sus componentes:
\begin{equation}
    \mu_p = \E[r_p] = \E\Bigl[\sum_{i=1}^{n} w_i r_i\Bigr]
    = \sum_{i=1}^{n} w_i \,\E[r_i]
    = \mathbf{w}^T \boldsymbol{\mu}.
    \label{eq:rentabilidad-cartera}
\end{equation}

La \textbf{varianza} de la cartera viene dada por la forma cuadrática
(\ref{eq:varianza-cartera}), que en notación matricial se escribe:
\begin{equation}
    \sigma_p^2 = \Var[r_p]
    = \mathbf{w}^T \boldsymbol{\Sigma} \mathbf{w}.
    \label{eq:varianza-cartera-matricial}
\end{equation}

\subsection{El problema de optimización media-varianza}

Markowitz formula la selección de carteras como un problema de
optimización bicriterio: maximizar $\mu_p$ y minimizar $\sigma_p^2$.
Esto da lugar a dos formulaciones equivalentes:

\begin{enumerate}
    \item \textbf{Minimizar el riesgo} dada una rentabilidad objetivo
          $\mu^*$:
          \begin{equation}
              \begin{aligned}
                  \min_{\mathbf{w}} \quad & \tfrac{1}{2} \mathbf{w}^T
                      \boldsymbol{\Sigma} \mathbf{w} \\
                  \text{s.a.} \quad & \mathbf{w}^T \boldsymbol{\mu}
                      \ge \mu^*, \\
                  & \sum_{i=1}^{n} w_i = 1.
              \end{aligned}
              \label{eq:qp-riesgo}
          \end{equation}

    \item \textbf{Maximizar la rentabilidad} dado un nivel de riesgo
          $\sigma_*$:
          \begin{equation}
              \begin{aligned}
                  \max_{\mathbf{w}} \quad & \mathbf{w}^T \boldsymbol{\mu} \\
                  \text{s.a.} \quad & \mathbf{w}^T \boldsymbol{\Sigma}
                      \mathbf{w} \le \sigma^2_*, \\
                  & \sum_{i=1}^{n} w_i = 1.
              \end{aligned}
              \label{eq:qp-rentabilidad}
          \end{equation}
\end{enumerate}

La primera es un \textbf{programa cuadrático convexo} (QP); la segunda, un
problema con \emph{restricción} cuadrática (QCQP), igualmente convexo. Ambas
producen la misma frontera eficiente (Capítulo~\ref{ch:formulacion}); se usa la
primera, más conveniente. Además del presupuesto ($\sum_i w_i = 1$), se imponen
dos restricciones por realismo: \textbf{long-only} ($w_i \ge 0$, sin ventas en
corto) y un \textbf{tope por activo} ($w_i \le w_{\max}$, típicamente $0{,}20$,
que limita la concentración).

%--------------------------------------------------------------------
\section{Frontera eficiente y cartera tangente}
%--------------------------------------------------------------------

\subsection{Definición de frontera eficiente}

Al resolver el problema (\ref{eq:qp-riesgo}) para distintos valores
del objetivo de rentabilidad $\mu^*$, se obtiene una familia de carteras
óptimas. En el
plano $(\sigma_p, \mu_p)$, estas carteras describen una curva denominada
\textbf{frontera eficiente}: para cada nivel de riesgo, la cartera que
maximiza la rentabilidad esperada; equivalentemente, para cada nivel de
rentabilidad, la cartera que minimiza el riesgo.

Las carteras por debajo de la frontera son \textbf{ineficientes} ---otra ofrece
más rentabilidad al mismo riesgo, o menos riesgo a igual rentabilidad--- y las
que quedan por encima son \textbf{inalcanzables} con el universo y las
restricciones dadas.

\subsection{Cartera de mínima varianza global (GMV)}

El extremo izquierdo de la frontera es la \textbf{cartera de mínima
varianza global} (GMV, por sus siglas en inglés). Se obtiene resolviendo
el QP sin restricción de rentabilidad objetivo:
\begin{equation}
    \begin{aligned}
        \min_{\mathbf{w}} \quad & \mathbf{w}^T \boldsymbol{\Sigma}
            \mathbf{w} \\
        \text{s.a.} \quad & \sum_{i=1}^{n} w_i = 1, \quad
            \mathbf{w} \ge \mathbf{0}.
    \end{aligned}
    \label{eq:gmv}
\end{equation}

La GMV es la cartera de menor riesgo posible dentro del universo de
activos y no depende de $\boldsymbol{\mu}$ (solo de $\boldsymbol{\Sigma}$).
Esta propiedad la convierte en una referencia especialmente útil: su
rendimiento depende exclusivamente de la calidad de la estimación de la
matriz de covarianza.

\subsection{Cartera tangente y máximo ratio de Sharpe}

El \textbf{ratio de Sharpe} \cite{Sharpe1966} mide el exceso de
rentabilidad por unidad de riesgo:
\begin{equation}
    S_p = \frac{\mu_p - r_f}{\sigma_p},
    \label{eq:sharpe}
\end{equation}
donde $r_f$ es la tasa libre de riesgo (que en este trabajo se toma como
$r_f = 0$ por simplicidad, dado que los tipos de interés de referencia se
mantuvieron próximos a cero durante buena parte del período de estudio).

La \textbf{cartera tangente} (o de máximo ratio de Sharpe) es el punto de la
frontera cuya recta desde el activo libre de riesgo es tangente a ella; según el
teorema de separación de Tobin \cite{Tobin1958}, es la combinación óptima de
activos con riesgo para cualquier inversor que preste o tome prestado al tipo
$r_f$. Con restricciones de desigualdad no admite forma cerrada y se obtiene
numéricamente, barriendo la frontera y tomando el punto de mayor Sharpe ---el
enfoque de nuestra implementación.

%--------------------------------------------------------------------
\section{Capital Market Line y CAPM}
%--------------------------------------------------------------------

\subsection{Capital Market Line (CML)}

En el plano $(\sigma_p, \mu_p)$, la \textbf{línea del mercado de capitales}
(CML) es la semirrecta que parte de $(0, r_f)$ y es tangente a la frontera
eficiente en la cartera tangente $T = (\sigma_T, \mu_T)$. Su ecuación es:
\begin{equation}
    \mu_p = r_f + \frac{\mu_T - r_f}{\sigma_T} \, \sigma_p.
    \label{eq:cml}
\end{equation}

Cualquier punto sobre la CML representa una combinación del activo libre
de riesgo y la cartera tangente $T$, y domina a cualquier cartera sobre
la frontera eficiente original (excepto en el punto de tangencia).

\subsection{Capital Asset Pricing Model (CAPM)}

El CAPM, desarrollado por Sharpe \cite{Sharpe1964} y Lintner
\cite{Lintner1965}, extiende el modelo de Markowitz introduciendo el
concepto de equilibrio de mercado. Su ecuación fundamental es:
\begin{equation}
    \E[r_i] - r_f = \beta_i \bigl(\E[r_m] - r_f\bigr),
    \label{eq:capm}
\end{equation}
donde $\beta_i = \Cov[r_i, r_m] / \Var[r_m]$ mide la sensibilidad del
activo $i$ a los movimientos del mercado, y $r_m$ es la rentabilidad
de la cartera de mercado.

El CAPM predice que, en equilibrio, toda la rentabilidad esperada de un
activo se explica por su exposición al riesgo sistemático (no
diversificable), medido por $\beta_i$. El riesgo idiosincrásico no se
remunera porque puede eliminarse diversificando.

Aunque el CAPM es un modelo de equilibrio y el de Markowitz uno normativo,
comparten la estructura media-varianza. Su conexión es relevante aquí en dos
sentidos: justifica usar el SPY (proxy de la cartera de mercado) como baseline,
y la estimación de $\boldsymbol{\Sigma}$ mediante PCA y factores
(Capítulo~\ref{ch:limitaciones}) responde a la misma idea de separar el riesgo
sistemático del idiosincrásico.

%--------------------------------------------------------------------
\section{Síntesis y conexión con el resto del trabajo}
%--------------------------------------------------------------------

El modelo de Markowitz proporciona un marco riguroso para la selección de
carteras, pero su implementación requiere estimar $\boldsymbol{\mu}$ y
$\boldsymbol{\Sigma}$, y la calidad de esas estimaciones determina el
rendimiento out-of-sample. De ello se ocupan los capítulos siguientes: el
Capítulo~\ref{ch:formulacion} desarrolla la formulación del QP y sus
condiciones de optimalidad; el Capítulo~\ref{ch:limitaciones}, el error de
estimación y los estimadores robustos de $\boldsymbol{\Sigma}$ (shrinkage,
PCA); y el Capítulo~\ref{ch:ml}, el uso de Machine Learning para estimar
$\boldsymbol{\mu}$.
